Fix an integer \(r\ge 2\). An integer \(n\ge 1\) is \emph{\(r\)-powerful} if for every prime \(p\mid n\) one has \(p^r\mid n\).

The problem asks:
(a) For \(r\ge 4\): does there exist a solution to
\[a_1+a_2+\cdots + a_{r-2}=c\]
where each \(a_i\) is \(r\)-powerful, the \(a_i\) are ``coprime'' (ambiguous; see below), and \(c\) is also \(r\)-powerful?
(b) If such solutions exist for \(r\ge 4\), are there only finitely many?
(c) For \(r=3\): are there infinitely many triples of coprime \(3\)-powerful numbers \((a,b,c)\) with \(a+b=c\)?

Ambiguity: For more than two integers, ``coprime'' may mean either (i) pairwise coprime, i.e. \(\gcd(a_i,a_j)=1\) for all \(i\ne j\), or (ii) collectively coprime, i.e. \(\gcd(a_1,\dots,a_{r-2})=1\). The example displayed in the problem statement for \(r=5\) satisfies (ii) but not (i).
